ALLES NICHT SEHR KLAUSURRELEVANT

Bruttoinlandsprodukt - Modellierung

Das BIP wird in der volkswirtschaftlichen Modellierung in der Regel durch den Buchstaben \( Y \) symbolisiert, es gilt:
\(
Y = C + I + G + EX - IM
\)

Dabei bezeichnen:
C := Konsumausgaben der privaten Haushalte
I := Investitionen
G := Öffentliche Ausgaben
EX := Kauf ausländischer Waren und Dienstleistungen durch Inländer
IM := Kauf inländischer Waren und Dienstleistungen durch Ausländer
S := Gesamtwirtschaftliche Ersparnis (hier gilt: \( S = Y - C \))

Es sei angenommen, dass es sich um eine geschlossene Volkswirtschaft handelt, d.h.: \(EX = IM = 0\)

damit gilt: \(Y = C + I + G\)

Angenommen, die Staatsausgaben sind ebenfalls 0 (\( G = 0 \)): \(Y = C + I \iff Y = Y - S + I \iff S = I\)

Die gesamtwirtschaftliche Ersparnisse entspricht der
gesamtwirtschaftlichen Investition

Reale Größen sind um die Effekte der Inflation bereinigt. Die Inflation ist wie folgt definiert:


Definieren Sie den VWL Begriff der Inflation
\( \pi_t = \frac{P_t - P_{t-1}}{P_t} \)

Die Inflation ist folglich eine Maßgröße für die Steigerung des Preisniveaus über einen gewissen Zeitraum (i.d.R. 12 Monate).
Die Inflation hat eine hohe Bedeutung im Rahmen der ökonomischen Betrachtung:
Die Messung der Inflation und des Preisniveaus erfolgt von staatlicher Seite anhand von definierter Warenkörbe

DER BMI (BIG MAC INDEX)

AD Kurve (Aggregated Demand)
• Die AD-Kurve beschreibt die gesamtwirtschaftliche Nachfrage.
• Die Ableitung erfolgt auf Basis des sog. IS-LM Modell.
Ein Punkt auf der Kurse repräsentiert eine \( (P,Y) \) -Kombination, für
die der Gütermarkt und der Geldmarkt im Gleichgewicht sind.
• AD-Kurve basiert damit auf der bekannten Gleichung
\( Y = C + I + G + EX – IM \)

INSERT IMAGE SLIDE 49

• Die negative Steigung werden insbesondere durch drei Effekte erklärt:
a) Vermögenseffekt
b) Zinseffekt
c) Wechselkurse.

INSERT IMAGE SLIDE 50


Vermögenseffekt
\( Y = C + I + G + EX – IM \)
Vermögenseffekt beschreibt Änderungen in der Nachfrage nach
Konsumgütern, die durch eine Erhöhung oder Reduktion an realem
Vermögen erzeugt werden.

AD Kurve –Zinseffekt
Zinseffekt
Der Zinseffekt beschreibt Änderungen in der Nachfrage nach
Investitionen, durch eine Änderung des Zinsniveaus.

Wechselkurseffekt
Der Zinseffekt beschreibt Änderungen in der Nachfrage nach
Investitionen, durch eine Änderung des Zinsniveaus.

Gesteigerte Konsumentennachfrage:
• Veränderung im Konsumentenverhalten
• Geringe Einkommensbesteuerung

DIFFERNZIEREN SIE ZWISCHEN KURVENVERSCHIEBUNG
ODER VERSCHIEBEN DES PUNKTES AUF DER KURVE 

Was ist die AS-Kurve und ihre Faktoren

(Aggregated Supply)
Die AS-Kurve beschreibt das aggregierte Angebot, das
Unternehmen bei unterschiedlichen Preisniveaus anzubieten bereit
sind. Das aggregierte Angebot wird durch die Grenzproduktivitäts-
theorie bestimmt. Einflussfaktoren sind dabei die Technologie
(Produktionsfunktion, der Kapitalbestand (Realkapital,
Humankapital)) sowie die Kosten (Preise und Löhne)
• Wesentliche Ursachen sind die Lohnsetzung und Preissetzung.


Setup
- Axes: real output (Y) on horizontal; price level (P) on vertical.
- Curves:
  - AD (Aggregate Demand): downward sloping. Captures planned spending: Y = C + I + G + NX.
  - SRAS (Short-Run Aggregate Supply): upward sloping due to sticky wages/prices and incomplete information.
  - LRAS (Long-Run Aggregate Supply): vertical at potential output Y* (full-employment output), which is assumed
  to be independent of prices in the long term (since the output mainly depends on the capacities and resources of a nation)

How curve shifts affect short-run equilibrium (intersection of AD and SRAS)
- AD shifts:
  - Rightward AD shift → higher Y and higher P (demand-pull expansion; unemployment falls).
  - Leftward AD shift → lower Y and lower P (recessionary pressure; unemployment rises).

- SRAS shifts:
  - Rightward SRAS shift → higher Y and lower P (positive supply shock; disinflation).
  - Leftward SRAS shift → lower Y and higher P (negative supply shock; stagflation).

- LRAS shifts (potential output changes):
  - Rightward LRAS → higher Y*; for a given AD, tends to lower P in the new long-run equilibrium.
  - Leftward LRAS → lower Y*; for a given AD, tends to raise P in the new long-run equilibrium.

Shifts of AD (anything that changes C, I, G, NX or the policy/financial conditions behind them)
- Monetary policy:
  - Rate cuts, QE, credit easing → AD right.
  - Rate hikes, QT, tighter credit → AD left.
- Fiscal policy:
  - Higher G or tax cuts/transfers that raise disposable income → AD right.
  - Lower G or tax hikes → AD left.
- Private sector conditions:
  - Higher consumer confidence/wealth; easier lending; asset booms → AD right.
  - Deleveraging, tighter lending, asset busts → AD left.
- External sector:
  - Foreign income boom → NX up → AD right.
  - Currency depreciation → NX up → AD right (but may raise import costs; see SRAS).
  - Currency appreciation or foreign recession → NX down → AD left.

Shifts of SRAS (unit production cost and expected price level)
- Input costs:
  - Oil/energy spike, higher import/input prices, currency depreciation raising import costs → SRAS left.
  - Energy price fall, cheaper inputs, currency appreciation lowering import costs → SRAS right.
- Wages and expectations:
  - Rising expected inflation or catch-up wage bargains → SRAS left.
  - Falling expected inflation or wage moderation → SRAS right.
- Productivity and supply-side frictions:
  - Productivity gains, deregulation reducing costs, lower payroll or production taxes → SRAS right.
  - Disruptions (natural disasters, supply-chain bottlenecks, new regulations raising costs) → SRAS left.

Shifts of LRAS (the economy’s capacity)
- Factors that raise capacity (LRAS right):
  - Capital deepening, technological progress, human capital, higher labor force participation/immigration, better institutions, infrastructure.
- Factors that lower capacity (LRAS left):
  - Capital destruction, persistent skill losses, adverse demographics, institutional breakdown, persistent scarring.



Bruttoinlandsprodukt - Modellierung

Das BIP wird in der volkswirtschaftlichen Modellierung in der Regel durch den Buchstaben \( Y \) symbolisiert, es gilt:
\(
Y = C + I + G + EX - IM
\)

Dabei bezeichnen:
C := Konsumausgaben der privaten Haushalte
I := Investitionen
G := Öffentliche Ausgaben
EX := Kauf ausländischer Waren und Dienstleistungen durch Inländer
IM := Kauf inländischer Waren und Dienstleistungen durch Ausländer
S := Gesamtwirtschaftliche Ersparnis (hier gilt: \( S = Y - C \))

Es sei angenommen, dass es sich um eine geschlossene Volkswirtschaft handelt, d.h.: \(EX = IM = 0\)

damit gilt: \(Y = C + I + G\)

Angenommen, die Staatsausgaben sind ebenfalls 0 (\( G = 0 \)): \(Y = C + I \iff Y = Y - S + I \iff S = I\)


INSERT IMAGE SLIDE 56

Eine bestimmte Produktionstechnologie wird durch die
Produktionsfunktion selbst ausgedrückt:
\( Y_s = f(K, L(\frac{w}{p})) \)

INSERT IMAGE SLIDE 59

FISKALPOLITIK
magisches viereck (Wirtschaftspolitik)
- Multipliktoreffekt

GELDPOLITIK

Barwert/Zins Anleihenbewertung



Skizieren sie die wesentlichen Wirtschafskreisläufe des Staates
Bruttoinlandsprodukt Kreislaufmodell
INSERT SLIDE 44

Akteure des Finanzsystems
INSERT SLIDE 45


Was sind Kernelemente der Wirtschaftspolitik?

Das magische Viereck:
- Außenwirtschaftliches Gleichgewicht
- Preisstabilität
- Wirtschaftswachstum
- Hoher Beschäftigungsgrad


Auf Basis des AS/AD Model kann abgeleitet werden, dass externe Einflüsse – z.B.:
durch Geldpolitik – Einfluss auf das realisierte Gleichgewicht der Gesamtökonomie hat
• Die externe Einflüsse können bspw. durch staatliches Handeln bestimmt werden.
Die Gesamtheit der Maßnahmen, mit denen der Staat in die regulierend in die
Wirtschaftssystem eingreift werden unter dem Begriff Wirtschaftspolitik zusammengefasst.
• In der ökonomischen Theorie wird unterschieden zwischen Fiskalpolitik und Geldpolitik
als staatliche Steuerungsobjekte
• Ziele der Wirtschaftspolitik werden im sog. magischen Viereck zusammengefasst.
Die hier verankerten wirtschaftspolitischen Ziele stehen dabei teilweise in Konkurrenz
zu einander.

Gegenstand der Fiskalpolitik ist der Versuch die auftretenden Konjunkturzyklen zu glätten.
Hierbei wird in Phasen eines Abschwungs/einer Rezession eine Expansive Fiskalpolitik
betrieben und in Fälle eines Aufschwungs/eines Booms eine kontraktive Fiskalpolitik
• Wichtige Instrumente der Fiskalpolitik sind die Staatsausgaben (z.B. Beauftragung von
privater Firmen durch den Staat), sowie die Höhe der Steuererhebungen
• Zusätzlich existieren weitere Stabilisationsmechanismen der sozialen Marktwirtschaft,
wie die Arbeitslosenversicherung

Im Gegensatz hierzu steht eine alternative Gestaltung der Fiskalpolitik:
angebotsorientierte Politik
- Gegenstand ist hier jeweils die Verbesserung der Angebotsbedingungen
z.B. durch Verbesserung der Wettbewerbsfähigkeit (u.a. Bürokratieabbau, Deregulierung)
sowie der Schaffung investitionsfreundlicher Rahmenbedingung

Ein wichtiger Wirkungsmechanismus in Bezug auf die Wirksamkeit der Fiskalpolitik
ist der sog. Multiplikatoreffekt:
• Der Staat erhöht seine Staatsausgaben, z.B. durch den Bau eines neuen
Regierungsgebäudes.
Eine Steigerung der Staatsausgaben erhöht die Wirtschaftsleistung überproportional.

Geldpolitik
• Unter dem Begriff der Geldpolitik werden die wirtschaftspolitischen Maßnahmen der
Zentralbank verstanden.
• Die Zentralbank der Eurozone ist die Europäische Zentralbank (EZB). Die wesentlichen
Aufgaben der EZB sind
- Sicherung der Preisstabilität (Mittelfristige Inflationsrate von etwa 2 Prozent)
- Unterstützung der wirtschaftspolitischen Ziele der EU, soweit diese nicht im Gegensatz zu
dem Ziel der Preisstabilität stehen.
- Aufgrund historischer Zeiten hoher Inflation ist heute ein zentrales Aufbauelement der
Wirtschaft, eine von der Politik unabhängige Zentralbank.
• Instrumente die der EZB zur Verfügung stehen um die Geldmenge und die Marktzinsen zu
beeinflussen:
- Offenmarktoperationen
- Änderung der Zentralbankzinssätze
- Änderung der Mindestreservepflicht

Unter der konventionellen Geldpolitik werden Maßnahmen zusammengefasst die mit Hilfe
des klassischen Instrumentariums der Notenbanken ergriffen werden.
Beispiele:
• Hauptrefinanzierungsgeschäfte: Die Zentralbank gewährt – gegen Gewährung
von Sicherheiten – Banken einen kurzfristigen Kredit zum sog. Hautrefinanzierungssatz
(„Leitzins“)
• Längerfristiges Refinanzierungsgeschäft: Es handelt sich hierbei ebenfalls um einen
gewährten Kredit. Dieser hat allerdings eine längere Laufzeit (z.B. über 12 Monate)
• Spitzenrefinanzierungsfazilität: Gewährung von Krediten über Nacht gegen Zinsen
• Einlagefazilität: Annahme von Einlagen über Nacht gegen Gewährung eines Zinses

Unkonventionelle Geldpolitik
Zielsetzung dieser Maßnahmen ist insbesondere eine Wirkung auf die langfristige Zinsstruktur.
Beispiele:
• Direkter Ankauf von Wertpapieren
• Der Zins für kurzfristige Papiere hatte bereits die Nullgrenze unterschritten, aber die
Anleihezinsen für länger laufende Papiere nicht
• Durch den Kauf der langlaufenden Papiere kann die Zinskurve hier abgeflacht werden


Stellen Sie klassische Ansichten der Volkswirtschaftslehre für Geldpolitik gegenüber
Keynsianismus
Grundlegende Annahmen:
• Der private Konsum ist grundsätzlich vom zur Verfügung stehenden
Einkommen abhängig
• Das Versagen von Märkten (z.B. durch träge Anpassungsprozesse)
führt zu Konjunkturschwankungen, Arbeitslosigkeit und Inflation
Ziel der Geldpolitik sind die kurzfristigen Zinsstrukturen

Wirtschaftspolitische Implikationen:
• Wirtschaftspolitik sollte antizyklisch betrieben werden (Fiskal- & und
Geldpolitik) und vorhandene Unsicherheiten reduzieren
Die Lohnpolitik ist eine wichtige Steuerungsgröße der
Konsumnachfrage

Monetarismus (Hayek, Friedmann)
Grundlegende Annahmen:
Der private Konsum ist abhängig vom permanenten Einkommen
(Vermögen) und damit weniger schwankunsanfällig durch kurzfristige
Marktgleichgewichte sind durch die natürliche Arbeitslosenrate
geprägt
Die Geldnachfrage ist relativ Stabil, Geldpolitik wird dazu genutzt das

Wirtschaftspolitische Implikationen:
• Geldpolitik wirkt lediglich kurzfristig, langfristig führt Geldpolitik nur zu
einer Anpassung des Preisniveaus (Inflation)
• Fiskalpolitik (staatlicher Konsum) führt dazu, dass der private Sektor
weniger konsumiert (Crowding Out)
Preisniveau stabil zu halten
Einkommensschwankungen
